\usepackage[utf8]{inputenc}
\usepackage[T1]{fontenc}
\usepackage{lmodern}
\usepackage{xspace}
\usepackage{xcolor}
\usepackage{mathtools}
\usepackage{bm}
\usepackage{relsize}
\usepackage{listings}
\usepackage{enumitem}
\usepackage{tikz}
\usepackage{manfnt}
\usepackage{wasysym}
\usepackage{array}
\usepackage{longtable}
\usepackage{fancyhdr}
\usepackage{lastpage}
\usepackage{xstring}
\usepackage{currfile}
\usepackage{xspace}
\usepackage{verbatim}
\usepackage[vscale=0.85,hscale=0.87]{geometry}
\usepackage{babel}
\usepackage{hyperref}


\graphicspath{{./images/}}

\edef\fichierPreambule{\currfilename}

% La version du paquet. Il faut que le fichier de documentation
% respecte la nomenclateur doc_v<version>.<extension>.
\def\version{\StrBehind*{\currfilebase}{doc_v}}


\setlength{\parindent}{0.8em}


\pagestyle{fancy}
\renewcommand{\headrulewidth}{0pt}
\fancyhead[LCR]{}
\fancyfoot[R]{\thepage/\pageref{LastPage}}
\fancyfoot[LC]{}

\hypersetup{%
  colorlinks,% colorie les liens
  urlcolor=cyan,% couleur des \url ou \href
  linkcolor=red,% couleur des liens internes (titres de section etc.)
}


\makeatletter
\renewcommand{\LT@makecaption}[3]{%
\LT@mcol\LT@cols c{\hbox to\z@{%
\hss\parbox[t]\LTcapwidth{\sbox\@tempboxa{#1{#2\CaptionSeparator}#3}%
\ifdim\wd\@tempboxa>\hsize #1{#2\CaptionSeparator}#3\else%
\hbox to\hsize{\hfil\box\@tempboxa\hfil}\fi\endgraf\vskip\baselineskip}\hss}}}
\makeatother 



% L'environnement alerte                        
\newsavebox{\boiteBrouillon}
\newcommand{\virageDanger}{\textdbend}

\newlength{\LargeurBouleAlerte}
\settowidth{\LargeurBouleAlerte}{%
  \begin{tikzpicture}%
    \node{\virageDanger};%
  \end{tikzpicture}%
}

% Style pour la boîte alerte
\tikzstyle{boitealerte}=[draw=red,rounded corners,inner xsep=1em,inner ysep=1ex]

% Style pour la boule alerte
\tikzstyle{boulealerte}=[circle,ball color=red,text=white] 

\newenvironment{alerte}{%
  \begin{lrbox}{\boiteBrouillon}% On sauve dans \boiteBrouillon le contenu
    \begin{minipage}{.8\linewidth}%  
      \color{red}%
      \setlength{\parskip}{1ex plus 0.2ex minus 0.2ex}%
}{%
    \end{minipage}%
  \end{lrbox}% Fin. Attention lrbox stocke du contenu sur 1 ligne (pas de paragrpahe)
  % La boîte peut être utilisée via \usebox{\boiteBrouillon}
  \vspace{1.75ex}% Si on met 1.5ex comme en bas, alors visuellement l'espace semble plus petit qu'en bas.
  \begin{tikzpicture}%
    \node [boitealerte] (cadre) {%
      \hspace{0.5\LargeurBouleAlerte}%
      \usebox{\boiteBrouillon}%
    };%
    \node [boulealerte] (alerte) at (cadre.west) {\virageDanger};%
  \end{tikzpicture}%
  \vspace{1.5ex}%
}   


% L'indentation des listes (ainsi que le réglage du label)
\setenumerate[1]{label=\arabic*.,ref=\arabic*,leftmargin=2em}
\setenumerate[2]{label=\alph*.,ref=\theenumi.\alph*,leftmargin=2em}
\setitemize[1]{label=\textbullet,leftmargin=2em}
\setitemize[2]{label=--,leftmargin=2em}

% Pour citer verbatim du texte comme les noms de menu etc.
% Les caractères spéciaux comme _, $ etc. doivent être échappés.
\newcommand{\verbtexte}[1]{%
  \begingroup%
  \bfseries%
  \color{teal}
  %\smaller%
  \texttt{#1}%
  \endgroup%
}

\lstset{language=bash,%
        basicstyle=\ttfamily,%
        commentstyle=\color{gray}\small,%
        showspaces=false,%
        showstringspaces=false,
        numbers=none,%left
        frame=single,%
        keepspaces=true,%
        flexiblecolumns=true,%
        morekeywords={elif},
        literate={é}{{\'e}}1 {è}{{\`e}}1 {ê}{{\^e}}1 {à}{{\`a}}1 %
                 {î}{{\^i}}1 {É}{{\'E}}1 {ç}{{\c c}}1 {ù}{{\`u}}1 %
                 {ï}{{\"i}}1 {â}{{\^a}}1 {û}{{\^u}}1,%
}

\newcommand{\return}[1]{\color{orange}\textbf{<Touche ENTR\'EE>}}

\newcommand{\Sedu}{SambaÉdu3\xspace}
\newcommand{\email}[1]{\href{mailto:#1}{\nolinkurl{#1}}}

\newenvironment{RQ}{%
  \begin{center}%
  \begin{minipage}{0.85\linewidth}%
  \smaller%
  \textbf{Remarque :}%
}{%
  \end{minipage}%
  \end{center}%
}


%% La macro \HeureCompilation donne l'heure de la compilation sous
%% la forme 21h37
\newcount\Heure
\newcount\Minute
\newcommand{\HeureCompilation}{%
  \Heure=\time %
  \divide \Heure by 60 %
  \multiply \Heure by 60 %
  \Minute=\time %
  \advance\Minute by -\Heure %
  \divide\Heure by 60 %
  \the\Heure h% ou bien heure\ifnum\Heure>1 s\fi
  \ifnum\Minute<10 %
     0\the\Minute %  ou bien minute\ifnum\Minute>1 s\fi
   \else %
     \the\Minute %
   \fi %
}


%\newcommand{\Vers}{\rightarrow}
\newcommand{\Vers}{$\rightarrow$\xspace}